\originalchapter{8.513作业1--4：几何动力学与玻色场}
\psource{Xiao-Gang Wen，8.513 Fall 2021，公开作业1--4；多数题取$\hbar=e=c=1$，需要时明确恢复。}
\originalsection{作业1}
\begin{exercise}[原题1(a)--(c)]
系统满足$\hbar b_{ij}\dot\xi^j=\partial_iH$，$b_{ij}=\partial_ia_j-\partial_ja_i$可逆，逆矩阵为$b^{ij}$。(a)证明$\hbar\dot O=[O,H]_{\rm cl}$，$[A,B]_{\rm cl}=b^{ij}\partial_iA\partial_jB$。(b)证明反对称性。(c)证明无显含时哈密顿量守恒。
\end{exercise}
\begin{solution}
左乘逆矩阵得$\hbar\dot\xi^i=b^{ij}\partial_jH$，链式法则即(a)。$b^T=-b$使$(b^{-1})^T=-b^{-1}$，交换$A,B$并重命名指标得负号，故$[H,H]_{\rm cl}=0$。若$H$显含$t$则$\dot H=\partial_tH$，不能省这个条件。Jacobi恒等式还来自$b$闭合，见第\ref{ch:berry}章的一般辛结构。
\end{solution}
\begin{exercise}[原题2(a)--(c)]
二维强磁场的导引中心满足$\dot x^i=B^{-1}\epsilon^{ij}\partial_jV$，$\epsilon^{12}=1$。(a)找相空间作用量。(b)量子化并给哈密顿算符。(c)若$V=v[(x^1)^2+(x^2)^2]/2$，求能级。这是最低Landau级的有效导引中心理论。
\end{exercise}
\begin{solution}
(a)取$L=-(B/2)(x\dot y-y\dot x)-V$，其$b_{xy}=-B$，逆$b^{xy}=1/B$，给题示方程。磁场与电荷的方向以所给方程固定，不另套相反导引中心符号。(b)$[\hat x,\hat y]=\ii\hbar/B$，$\hat H=V(\hat x,\hat y)$，一般非二次$V$需说明算符排序，此题无歧义。(c)取$\hat Q=\sqrt{|B|}\hat x,\hat P=\operatorname{sgn}(B)\sqrt{|B|}\hat y$，$[Q,P]=\ii\hbar$，$H=v(Q^2+P^2)/(2|B|)$，所以$E_n=\hbar v(n+1/2)/|B|$，$v>0$时有界。被投影掉的回旋零点能是共同常数，不含在$V$的有效谱中。
\end{solution}
\begin{exercise}[原题3(a)--(b)]
(a)$a=a_x\dd x+a_y\dd y$，$b=\dd a=(\partial_xa_y-\partial_ya_x)\dd x\wedge\dd y$，证明圆盘的Stokes公式。(b)在三维原点外，$b=r^{-3}(x\dd y\wedge\dd z-y\dd x\wedge\dd z+z\dd x\wedge\dd y)$，证明$\dd b=0$并求半径$r$球面的积分。
\end{exercise}
\begin{solution}
(a)这是平面Green公式：分别把$\partial_xa_y$沿$x$积分、$-\partial_ya_x$沿$y$积分，区域相邻切片的内边界相消，只剩逆时针定向外边界$\oint a$。先对矩形证明，再将光滑圆盘分割逼近，要求$a$及一阶导数足够连续。(b)三项对应向量场$\vct r/r^3$的通量二形式，其散度$3r^{-3}-3r^{-5}r^2=0$于$r\ne0$。球面拉回为$\sin\theta\dd\theta\wedge\dd\phi$，积分$4\pi$。它闭合但在穿孔空间不全局恰当：若$b=\dd a$全局光滑，闭球面积分应为0，与$4\pi$矛盾。
\end{solution}
\begin{exercise}[原题4(a)--(b)]
量子自旋$S$有$H=-\vct B\cdot\vct S$，$\vct B=B\hat z$。取沿单位向量$n$的最大本征值相干态。(a)求$n$的相空间作用量。(b)求经典运动。
\end{exercise}
\begin{solution}
自旋$1/2$相干态$(\cos(\theta/2),\e^{\ii\phi}\sin(\theta/2))^T$的$\ii\avg{n|\partial_tn}=-(1-\cos\theta)\dot\phi/2$。自旋$S$可视为$2S$个对称自旋$1/2$的相干乘积，Berry项相加，$L=-\hbar S(1-\cos\theta)\dot\phi+BS\cos\theta$。变分给$\dot\theta=0,\dot\phi=-B/\hbar$，即$\dot n=n\times B/\hbar$。另一规范可加$\hbar S\dot\phi$总导数，却不改变局部方程；球面曲率通量$-4\pi S$需用两个相位片，量子条件$2S\in\mathbb Z$保证全局一致。
\end{solution}
\originalsection{作业2}
\begin{exercise}[原题1(a)--(c)]
带电粒子在$L_x\times L_y$环面，$H=(-\ii\nabla-A)^2/(2m)$，$A=(-By,0)$，$\hbar=e=c=1$。(a)$\psi'=\e^{\ii f}\psi$时求$H'$。(b)保留$x$周期，求使$A(y+L_y)$与$A(y)$一致的规范函数及$y$边界；证明总磁通量子化。(c)证明最低Landau级简并等于通量数，并说明与自旋球面“态数比Chern数多1”的区别。
\end{exercise}
\begin{solution}
(a)$A'=A+\nabla f$，$H'=\e^{\ii f}H\e^{-\ii f}$。(b)$f=-BL_yx$，$U=\e^{-\ii BL_yx}$，所以$\psi(x,y+L_y)=\e^{-\ii BL_yx}\psi(x,y)$。先绕$x$再绕$y$与反次序应同态，要求$\e^{-\ii BL_xL_y}=1$，即$BL_xL_y=2\pi n$。这正是线丛过渡函数一致性。

(c)Landau规范导引中心随$x$动量$2\pi l/L_x$移动，$y$方向周期识别$l\sim l+n$，有$|n|$个独立中心。若$B=0$则不是Landau投影问题。球面自旋相空间曲率通量绝对Chern数$2S$，量子自旋Hilbert维数$2S+1$；不同拓扑与量子化的端点修正使环面的等号不能照搬到球面。
\end{solution}
\begin{exercise}[原题2(a)--(d)]
三角格的四子格磁矩为$n_1=(1,1,1)/\sqrt3,n_2=(1,-1,-1)/\sqrt3,n_3=(-1,1,-1)/\sqrt3,n_4=(-1,-1,1)/\sqrt3$。以三角原矢$a_1,a_2$，格点$(n,m)$的类型按偶偶、奇偶、偶奇、奇奇为1、2、3、4。电子自旋始终平行局域矩，$\ket{n_i}=(\cos\theta_i/2,\e^{\ii\phi_i}\sin\theta_i/2)^T$，$t_{ij}=A\avg{n_i|n_j},A>0$。(a)求沿逆时针三角环的均匀磁通。(b)重新选自旋态相位，使跃迁仅使胞加倍。(c)求两带和隙。(d)关掉$a_1$水平键后求新带与隙。条件由原图转为向量与格点标签，不复用原示意图。
\end{exercise}
\begin{solution}
(a)任意不同两方向内积为$-1/3$，故$|t|=A/\sqrt3$。投影矩阵$P_i=(1+n_i\cdot\sigma)/2$给Bargmann积$\Tr(P_iP_jP_k)=[1+n_i\cdot n_j+n_j\cdot n_k+n_k\cdot n_i+\ii n_i\cdot(n_j\times n_k)]/4$。前三内积使实部为0，按原图逆时针$1\to2\to3$，虚部正，通量$\pi/2$模$2\pi$；相邻反向排布补偿空间方向，所有逆时针三角同号。

(b)原相干态规范下相位依四子格，选择附加相位$(\zeta_1,\zeta_2,\zeta_3,\zeta_4)=(0,\pi/4,-\pi/4,-\pi/2)$。令$t=A/\sqrt3$，沿$a_1$为$t$，沿$a_2$为$(-1)^nt$，沿$a_1-a_2$为$-\ii(-1)^nt$，都与$m$无关且只依$n$奇偶，胞仅沿$a_1$倍增。

(c)以$k_i=k\cdot a_i$、保留两子格内部位置相位的Bloch规范，矩阵可取$H=2t[\cos k_1\sigma_x+\cos(k_1-k_2)\sigma_y+\cos k_2\sigma_z]$，$E_\pm=\pm2t\sqrt{\cos^2k_1+\cos^2k_2+\cos^2(k_1-k_2)}$。最小平方和为$3/4$，可由令$k_1=-k_2=\pi/3$达到，并由三角恒等式或求驻点与边界核对全局最小；全隙$2\sqrt3t=2A$。

(d)去掉$\sigma_x$项，$E_\pm=\pm2t\sqrt{\cos^2k_2+\cos^2(k_1-k_2)}$，在$k_2=\pi/2,k_1=0$等点闭隙，故隙0。少一类键改变的是完整二维Bloch矩阵，不能只把原带宽乘$2/3$。
\end{solution}
\originalsection{作业3}
\begin{exercise}[原题1(a)--(c)，斜入射选做]
自旋$1/2$粒子$H=p^2/(2m)+g\sigma\cdot p+V(x)$，$g>0$。(a)$V=0$求给定$p$最低能与自旋态。(b)该低态束沿$+z$、$p>mg$撞击$z>0$无限墙，反射保持自旋，求反射能、动量与速度。(c)取以$(x,p)$为中心、宽$\xi$的Gaussian波包，$V$在$\xi$内缓变，求相空间作用量与方程。斜入射可选做，须保留墙平行动量。
\end{exercise}
\begin{solution}
(a)$\sigma\cdot\hat p$本征值$\pm1$，低支$\epsilon_-=p^2/(2m)-g|p|$，自旋反平行$p$。(b)初自旋沿$-z$，固定该自旋的能量为$p_z^2/(2m)-gp_z$。等能第二根$p'_z=2mg-p$，反射速度$v'=p'/m-g=-(p/m-g)$，能量不变。若$p>2mg$，$p'<0$而自旋仍负，反射落在该动量的高螺旋支；不能强迫反射后仍为最低支而违反边界自旋守恒。

(c)在孤立低支的绝热条件下$L=p\cdot\dot x+\hbar\mathcal A_p\cdot\dot p-\epsilon_-(p)-V(x)$，$\Omega_p=\hat p/(2|p|^2)$（本册$\mathcal A=\ii\avg{u|\nabla u}$）。$\dot p=-\nabla V$、$\dot x=p/m-g\hat p-\hbar\dot p\times\Omega_p$。$p=0$两支闭隙，不满足单支波包条件。对硬墙本身也不能以缓变单支方程代替(b)的两支边界散射；斜入射须保持$p_x,p_y$并叠加允许的反射自旋本征支来满足$\psi(z=0)=0$，具体地，保持平行波矢$q=p\sin\theta$，固定能量$E=p^2/(2m)-gp$。可用壳半径为$r_1=p$及$r_2=|p-2mg|$；$p>2mg$时第二壳属于高螺旋支，$mg<p<2mg$时属于低支内壳。每壳取满足向外法向群速度的$p_{z,j}=\pm\sqrt{r_j^2-q^2}$；若根虚，选$\operatorname{Im}p_z<0$使波在$z\to-\infty$衰减。以$(g\sigma\cdot p_j)u_j=[E-r_j^2/(2m)]u_j$求对应两分量自旋矢，边界系数由$[u_1,u_2](A_1,A_2)^T=-u_{\rm in}$解出。传播支概率$R_j=|A_j|^2|v_{z,j}|/v_{z,\rm in}$，倏逝支不带远处净流。这个两根加边界的构造覆盖斜入射，法向$q=0$退化为(b)的自旋守恒单支反射；壳退化点需取连续极限。
\end{solution}
\begin{exercise}[原题2(a)--(b)]
石墨烯两谷近似$H_K=\hbar v(k_x\sigma_x+k_y\sigma_y)$、$H_{K'}=\hbar v(k_x\sigma_y+k_y\sigma_x)$，$\mu=0$。(a)求热激发导带电子及价带空穴密度的温度幂律。(b)假设$\tau$不随$T$，由Boltzmann求电导，判断是否与密度同为$T^2$。
\end{exercise}
\begin{solution}
每谷每自旋$D(E)=|E|/(2\pi\hbar^2v^2)$。设总简并$g_d$，$n_e=n_h=g_d(\kb T)^2\pi/(24\hbar^2v^2)$，由$\int_0^\infty x/(\e^x+1)\dd x=\pi^2/12$。电导$\sigma=e^2\tau(v^2/2)\int D_{\mathrm{total}}(E)[-f'(E)]\dd E$，正负带均计，$\int_0^\infty E[-f']\dd E=\kb T\log2$，得$\sigma=g_de^2\tau\kb T\log2/(2\pi\hbar^2)$，仅线性$T$。原因是Dirac载流子的惯性尺度$|E|/v^2$也随温度改变，不能固定质量套Drude密度比例。$T\to0$的真实量子极限另含无序、接触与带间效应。
\end{solution}
\begin{exercise}[原题3(a)--(c)]
在上述两谷都加$\Delta\sigma_z$。(a)求两带Berry曲率峰、每谷总通量与宽$\delta k$。(b)零温K谷有导带电子密度$n_e^K$，K'谷有价带空穴密度$n_h^{K'}$，求小、大密度的Hall响应并说明连续依赖。(c)K谷导带电子初始$x=k=0$，受$E=E\hat x$，短时间$\tau$内曲率近常数，求位置。
\end{exercise}
\begin{solution}
令$\eta_K=+1,\eta_{K'}=-1$，$E_k=\sqrt{\Delta^2+\hbar^2v^2k^2}$，本册约定$\Omega_{\eta,\pm}=\mp\eta\Delta\hbar^2v^2/(2E_k^3)$。峰宽$|\Delta|/(\hbar v)$，每谷单带全平面积分$\Phi_\pm=\mp\eta\pi\operatorname{sgn}\Delta$，完整晶格两谷合计为整数且此同号质量使总Chern为0。

设各密度包含自旋简并$g_s$，$k_F^2=4\pi n/g_s$。定义$F(n)=\frac{\operatorname{sgn}\Delta}{2}[1-|\Delta|/\sqrt{\Delta^2+\hbar^2v^2k_F^2}]$。满价带的两个谷响应相消，新增电子和移除空穴后的曲率积分为$g_s[-F(n_e^K)+F(n_h^{K'})]$。用$\sigma_{xy}=-e^2\int f\Omega/\hbar$的方向约定，$\sigma_{xy}=g_s(e^2/h)[F(n_e^K)-F(n_h^{K'})]$。小密度$F\simeq\operatorname{sgn}\Delta\,\pi\hbar^2v^2n/(g_s\Delta^2)$，大密度各自趋$\operatorname{sgn}\Delta/2$。相同谷相反带的电子--空穴贡献相加，而本题相反谷会抵消；等密度时净Hall为0，不能遗漏K'手性负号。

(c)$k_x=-eEt/\hbar$，带底群速度$x$分量$\hbar v^2k_x/|\Delta|$，因此$x(\tau)=-eEv^2\tau^2/(2|\Delta|)$。异常速度$\dot y=-eE\Omega_{K,+}(0)/\hbar$，所以$y(\tau)=-eE\Omega_{K,+}(0)\tau/\hbar$。适用$eEv\tau\ll|\Delta|$且跨带跃迁小；速度与外力的线性关系包含横向几何项。
\end{solution}
\originalsection{作业4}
\begin{exercise}[原题1(a)--(c)]
一维线性色散玻色场在相距$L$的两固定墙间，$k_n=n\pi/L$，零点能$\frac12\sum\hbar ck_n$。(a)热核正规化乘$\e^{-ak_n}$，求小$a$展开并忽略线性及更高项。(b)格点正规化$E_{\rm latt}=\sum_{n=1}^{N}c\hbar a^{-1}\sin(k_na/2),N=L/a$，同样展开。(c)三墙总长$L$，中墙距左墙$l$，求总能及$l\ll L$时力，区分共同紫外项与有限力。
\end{exercise}
\begin{solution}
(a)几何级数导数$\sum n q^n=q/(1-q)^2$，$x=\pi a/L$给$E_{\mathrm{heat}}=\pi c\hbar\e^{-x}/[2L(1-\e^{-x})^2]=c\hbar L/(2\pi a^2)-\pi c\hbar/(24L)+O(a^2/L^3)$。

(b)有限正弦和为$\sum_{n=1}^N\sin(n\pi/2N)=[\cot(\pi/4N)+1]/2$，所以$E_{\mathrm{latt}}=2c\hbar L/(\pi a^2)+c\hbar/(2a)-\pi c\hbar/(24L)+O(a^2/L^3)$。两方案体积与边界发散不同，有限$1/L$项相同。形式$\sum n=-1/12$在此是扣掉规定发散后的解析正规化值，不是普通正项级数的和。

(c)分别加$E(l)+E(L-l)$，共同体积项只依总$L$，格点边界项只依墙数，故对中墙力仅有限项贡献：$F_l=-\pi c\hbar[1/l^2-1/(L-l)^2]/24$，$l\ll L$时$F=-\pi c\hbar/(24l^2)$，负号指向近墙。移动墙时保持同一微观截止，不用两个随长度变化的截止“消掉”发散。
\end{solution}
\begin{exercise}[原题2(a)--(e)]
格点玻色模型$H=\sum_i[2t\phi_i^\dagger\phi_i-t(\phi_i^\dagger\phi_{i+1}+\mathrm{h.c.})-\mu\phi_i^\dagger\phi_i+U\phi_i^{\dagger2}\phi_i^2]$，$t,U>0$。(a)证明未归一相干态重叠$\langle0|\e^{\sum\phi_i^*\hat\phi_i}\e^{\sum\widetilde\phi_i\hat\phi_i^\dagger}|0\rangle=\e^{\sum\phi_i^*\widetilde\phi_i}$。(b)求归一及共同本征值。(c)求$L=\langle\phi|\ii\partial_t-H|\phi\rangle$。(d)$\mu\ge0$时求均匀极小。(e)线性化运动求$\omega_k$，比较$\mu>0,=0,<0$的低能与动态指数。
\end{exercise}
\begin{solution}
(a)湮灭指数移过创建指数产生数对易子$\sum\phi_i^*\widetilde\phi_i$，余指数对真空作用为1。(b)归一因子$\e^{\sum|\phi_i|^2/2}$，且湮灭算符可移过指数得本征值$\phi_i$。(c)$\bar H=\sum[t|\phi_{i+1}-\phi_i|^2-\mu|\phi_i|^2+U|\phi_i|^4]$，Berry项$\ii\sum\phi_i^*\dot\phi_i$差一总导数后即原题形式。

(d)均匀跃迁项0，极小$|\bar\phi|^2=\mu/(2U)$，相位任意，$\mu=0$在原点。(e)取$\bar\phi$实，$\ii\dot{\delta\phi_i}=t(2\delta\phi_i-\delta\phi_{i+1}-\delta\phi_{i-1})+\mu(\delta\phi_i+\delta\phi_i^*)$，一纯复平面波通常不闭合，需$u\e^{\ii(ki-\omega t)}+v^*\e^{-\ii(ki-\omega t)}$的Bogoliubov双分量。$\epsilon_k=2t(1-\cos k)$，$\omega_k=\sqrt{\epsilon_k(\epsilon_k+2\mu)}$。$\mu>0$长波$\sqrt{2t\mu}|k|$，$z=1$；$\mu=0$为$tk^2$，$z=2$；$\mu<0$均匀真空稳定，$\omega_k=\epsilon_k+|\mu|$有隙。该相干态平均场刻画真空--超流入口，不证明一维存在通常长程凝聚，相位涨落仍需第\ref{ch:boson}章处理。
\end{solution}
